% Unified upper bound slide 2: Bound the post-restart and restart accumulation terms.
% Replaces upper_term1.tex + upper_term2.tex (no equation labels).

\small
\textbf{1. Post-restart evolution.}
On $E_N$, $\tau\le T$ and the process after $\tau$ is a translated $\mathcal{V}^{N|\alpha}$:
\[
    \mathbb{E}\!\left[
        \bigl(\mathcal{M}_{\ell_N}^{\mathrm{b}}-\max\mathcal{V}^{\mathrm{b}}_\tau\bigr)\mathbf{1}_{E_N}
    \right]
    =
    \mathbb{E}\!\left[
        \bigl.\mathbb{E}\!\left[\mathcal{M}^{N|\alpha}_{\ell_N-\iota}\right]\bigr|_{\iota=\tau}\mathbf{1}_{E_N}
    \right]
    \le
    \left(-\alpha+\frac{\pi^2}{2}\sigma^2 t\right)\log N+o_N(\log N),
\]
where the last bound uses \textbf{Theorem~1.1} (the $\mathrm{L}^2$ limit theorem) for $a=\alpha\in(0,1)$.

\vspace{3mm}
\textbf{2. Restart accumulation.}
Define the step sums during restarts:
\[
    G_j:=\tau_j-\tau_{j-1}-1
    \;\sim\;
    \mathrm{Geom}(p_0),
    \qquad
    X_k:=\sum_{i=\tau_{k-1}+1}^{\tau_k-1}
    \bigl(\mathcal{V}^{\mathrm b}(w_{i+1})-\mathcal{V}^{\mathrm b}(w_i)\bigr),
    \qquad
    S_k:=\sum_{i=1}^k X_i,
\]
where $G_j$ (i.i.d.) is the \textbf{time spent} in the $j$-th reboot (until a two-child event occurs),
and $X_k$ is the \textbf{total displacement} during the $k$-th reboot
(a sum of $G_k$ i.i.d.\ copies of $\Xi_1$, the leftmost step).

From the construction, $\max\mathcal{V}^{\mathrm{b}}_\tau\le S_\kappa+(\kappa-1)M+\ell M$, and on $E_N$,
\[
    \mathbb{E}\!\left[\max\mathcal{V}^{\mathrm{b}}_\tau\mathbf{1}_{E_N}\right]
    \le
    \mathbb{E}[S_\kappa]+O_N(\log N).
\]

Since $\kappa$ is a stopping time w.r.t.\ i.i.d.\ step sums, \textbf{Wald's identity} gives
\[
    \mathbb{E}[S_\kappa]
    =
    \mathbb{E}[\kappa]\,\mathbb{E}[X_1]
    =
    \mathbb{E}[\kappa]\,\mathbb{E}[G_1]\,\mathbb{E}[\Xi_1]
    \le
    \frac{1}{p_1}\cdot\frac{1-p_0}{p_0}\cdot\bigl(\mathbb{E}[(\Xi_1)_+]+\mathbb{E}[(\Xi_1)_-]\bigr)
    =O_N(1).
\]

Thus the restart accumulation term is \textbf{bounded}, contributing only $O_N(\log N)$.
